\clearpage\thispagestyle{empty}
\pdfbookmark[0]{Quellen für Erste Beobachtungen von Le Verriers Planeten}{leverrierfirst-sources}
\begin{center}\textbf{Quellen für Erste Beobachtungen von Le Verriers Planeten}\end{center}
\noindent Diese KI-Übertragung ins heutige Deutsch folgt dem Bericht auf S.~471 der \href{https://gdz.sub.uni-goettingen.de/id/PPN236007467}{\textit{Werke}, Band~VI} (Göttingen, 1874). Heinrich Christian Schumacher gibt drei Beobachtungen von Carl Friedrich Gauß wieder. Der Erstdruck steht in der Beilage zu den \textit{Astronomischen Nachrichten}, Band~XXV, Nr.~581, \href{https://api.digitale-sammlungen.de/iiif/image/v2/bsb10049450_00049/full/1800,/0/default.jpg}{Sp.~82}. Die Quellenangabe der Werkausgabe nennt den 22. Oktober 1846. Der Katalog der Bayerischen Staatsbibliothek datiert den gebundenen Band auf 1847.
\par\smallskip\noindent Der Zeitschriftenbericht enthält auch Beobachtungen und Messungen anderer Astronomen. Die Werkausgabe gibt die drei Zeilen von Gauß mit den wiederholten Minuszeichen und den leeren Gradangaben wieder. Der Auszug endet mit Schumachers Unterschrift. KI-Übertragung und KI-Satz.
\clearpage
